\section{The actual sheaf lifts and the geometric normal
quotient}\label{the-actual-sheaf-lifts-and-the-geometric-normal-quotient}

Complete derivation, 24 September 2026. Stable proof locators CW0--CW11.

\subsection{CW0. Construction stage and exact
question}\label{cw0.-construction-stage-and-exact-question}

The support is \(\tau\langle Z_1;\mathrm{no}\ Z_2\rangle\). All
operations below are on the arithmetic coefficient sheaves already
constructed in CSP, after the complete arithmetic reconstruction. There
is no addition, scalar multiplication, coordinate, distance, or parity
operation on \(\tau\). The two charts remain separately labelled; a
difference in a cochain complex is not a pooled branch return measure.

The connected argument and its correction-precedence table in
\nolinkurl{USER_ARGUMENT_RECONSTRUCTION.md,} the operative
\nolinkurl{READ_FIRST_USER_CONSTRUCTION.md,} and the full governing
passages WU050--WU055 and WU061--WU063 in
\nolinkurl{USER_CONSTRUCTION_FULL_LOGBOOK.md} were read for this
calculation. In particular, no arithmetic label initializes the
supporting construction; the original zeta arithmetic is retained; and a
calculation on this receiving sheaf is not promoted into a conclusion
about every possible geometric realization of the complete construction.

The actual input is the final
\nolinkurl{CC_SPHERE_PULLBACK_AND_NORMAL_DIRECTION.md,} CSP0--CSP13,
read in full. Its version for this derivation has SHA256 \[
\texttt{EBFF9973E1472EA88E3F88D3350FAD8FFA98364067F485F48D3AF4E119EBAF9D}.
\] The two comparison proofs are
\nolinkurl{GLOBAL_SHIFTED_ZETA_LIFT.md,} GSL0--GSL11, and
\nolinkurl{DELIGNE_INVARIANT_CYCLE_QUOTIENT.md,} DC0--DC12, previously
derived and read in full and checked against their unchanged final
hashes \[
\texttt{2CC092DD3EBC0F021E9914FC7FD75DA16EA4DE47760B1D12558F33F0A661F529},
\] \[
\texttt{3A9AFDAEACD7DB0F154348F6487A204852459B08AE8E0F90017E7755FD627521}.
\] The action and inverse calculations used from those notes are proved
again below where they determine the present conclusion.

The source construction is Connes and Consani, \emph{Schemes over
\(\mathbb F_1\) and zeta functions},
\href{https://arxiv.org/abs/0903.2024v3}{arXiv:0903.2024v3}, §5: the
original author file
\nolinkurl{sources/CC_0903_2024_v3/author_source/announc3.tex,}
lines1376--1667, was read in the preceding CSL verification. The
relevant source formulas are fonction3, restmaps, bord, hzero, rep1,
liftw and court. The finite-unit-sector comparison retains CS2A's
full-difference generating-domain qualification. The later signed source
CC.tex, lines494--610 and773--899, was also read: its topological
comparison does not by itself identify its structure sheaf with this
complex coefficient sheaf.

The comparison with Deligne concerns \emph{La conjecture de Weil. II},
\href{https://www.numdam.org/item/PMIHES_1980__52__137_0/}{§3.6}. DC
reconstructs that argument from the current French transcription, with
its actual geometric hypotheses and arithmetic reduction. No new
original-author-TeX reading of Deligne is claimed here.

This note computes both actual sheaf short exact sequences in CSP, their
connecting maps, and the actual supported localization sequence. It
proves the vanishing and existence of the lifts that those sequences
define. It also computes exactly which kernel has a separated
original-zero character range, and which kernel does not.

\subsection{CW1. Full coefficients, original zeta, and the source
image}\label{cw1.-full-coefficients-original-zeta-and-the-source-image}

Use \[
S=\{f\in\mathcal S(\mathbb R;\mathbb C):
f(-v)=f(v),\ f(0)=0,\ \int_{\mathbb R}f(v)\,dv=0\},
\qquad
V_\pm=S\oplus E_\pm,\quad E_\pm=\mathbb C^2.
\tag{CW1.1}
\] The two endpoint coordinates are labelled value and integral on each
chart. Let \[
A=\left\{a\in C^\infty(\mathbb R_{>0}):
p_{N,j}(a)=\sup_{u>0}(u^N+u^{-N})
|(u\partial_u)^j a(u)|<\infty\quad(N,j\ge0)\right\}.
\tag{CW1.2}
\] The actual maps are \[
\Sigma f(u)=2\sum_{k\ge1}f(ku),\qquad
Ra(u)=u^{-1}a(u^{-1}),\qquad
r_+(f,c_0,c_1)=\Sigma f,\qquad
r_-(h,d_0,d_1)=R\Sigma h=\Sigma\widehat h.
\tag{CW1.3}
\] The factor \(2\) records both rational signs. The Fourier convention
is \(\widehat h(t)=\int_{\mathbb R}h(v)e^{-2\pi ivt}\,dv\). Before
imposing the two endpoint conditions, the retained Poisson identity is
\[
\Sigma\widehat h(u)
=u^{-1}\Sigma h(u^{-1})+u^{-1}h(0)-\int_{\mathbb R}h(v)\,dv.
\tag{CW1.4}
\] Thus \(R^2=1\), Fourier transformation preserves \(S\) and squares to
one there, and \[
J=\Sigma(S)=r_+(V_+)=r_-(V_-),\qquad
\ker r_\pm=E_\pm.
\tag{CW1.5}
\] Continuity follows by Schwartz summation for \(u\ge1\) and (CW1.4)
for \(u\le1\). Injectivity follows by taking the Mellin transform in
\(\Re s>1\), dividing by the nonzero Euler product there, and applying
Fourier uniqueness on a vertical line. These are the source maps of
CSP1, not new restrictions.

Write \[
\mathcal B=\{F\text{ entire}:
b_{M,N}(F)=\sup_{|\Re s|\le M}(1+|\Im s|)^N|F(s)|<\infty
\text{ for all }M,N\ge0\}.
\] The exact comparison is \[
\Theta a(s)=\frac12\int_0^\infty a(u)u^s\,\frac{du}{u},
\qquad
(\Theta^{-1}F)(u)=\frac{u^{-1/2}}{\pi}
\int_{\mathbb R}F(1/2+it)u^{-it}\,dt.
\tag{CW1.6}
\] The integrals and their differentiated versions converge by the
seminorms. Integration by parts in \(\log u\), followed by the inverse
Fourier integral and contour shifts in each fixed strip, proves that
\(\Theta:A\to\mathcal B\) is a topological isomorphism. The \(1/2\),
\(1/\pi\), and \(u^{-1/2}\) are retained in these comparison maps.

For \(\Re s>1\), absolute summation gives \[
\Theta\Sigma f(s)=\zeta(s)a_f(s),\qquad
a_f(s)=\int_0^\infty f(v)v^s\,\frac{dv}{v}.
\tag{CW1.7}
\] The original arithmetic remains \[
\zeta(s)=1+\sum_{k\ge2}k^{-s}
=\prod_p(1-p^{-s})^{-1},\qquad
-\frac{\zeta'(s)}{\zeta(s)}
=\sum_p\sum_{m\ge1}(\log p)p^{-ms}\quad(\Re s>1).
\tag{CW1.8}
\] In particular, neither the unit summand nor any prime-power
repetition is removed.

For the actual source \[
f_0(v)=\frac{\pi}{2}v^2(2\pi v^2-3)e^{-\pi v^2}
\] the full transform is \[
F_0(s)=\Theta\Sigma f_0(s)
=\frac{s(s-1)}8\pi^{-s/2}\Gamma(s/2)\zeta(s),
\qquad
\mathcal M_0\Sigma f_0=2F_0.
\tag{CW1.9}
\] Here \(f_0(0)=0\) and its integral is zero because \[
\frac{\pi}{2}
\left(2\pi\frac{3}{4\pi^2}-3\frac1{2\pi}\right)=0.
\] Its two Gaussian Mellin integrals, before applying the Gamma
recurrence, are \[
\frac{\pi}{2}\left[
2\pi\cdot\frac12\pi^{-(s+4)/2}\Gamma((s+4)/2)
-3\cdot\frac12\pi^{-(s+2)/2}\Gamma((s+2)/2)\right].
\tag{CW1.10}
\] Multiplication by the original \(\zeta(s)\) gives (CW1.9). The
endpoint and trivial-zero cancellations have the exact values \[
F_0(0)=F_0(1)=\frac18,
\] \[
F_0(-2r)
=\frac{r(2r+1)(-1)^r\pi^r}{2\,r!}\zeta'(-2r)
=\frac{(1+2r)(2r)}8\pi^{-(1+2r)/2}
\Gamma((1+2r)/2)\zeta(1+2r)\ne0\quad(r\ge1).
\tag{CW1.11}
\] The original \(\zeta\) still has its pole at \(1\) and its trivial
zeros at \(-2r\). For general \(f\in S\), the retained receiver values
are \[
\Theta\Sigma f(-2r)
=\zeta'(-2r)\frac{f^{(2r)}(0)}{(2r)!},\qquad
\Theta\Sigma f(1)=a_f'(1),\qquad
\Theta\Sigma f(0)=-\tfrac12a_f(0).
\tag{CW1.12}
\] Taylor subtraction in the integral defining \(a_f\) proves its
meromorphic continuation. At \(\nu=-2r\), write \[
a_f(\nu+t)=A_{-1}/t+\sum_{k\ge0}A_kt^k,\qquad
A_{-1}=f^{(2r)}(0)/(2r)!.
\] Then the coefficient of \(t^n\) in \(\zeta(\nu+t)a_f(\nu+t)\) is \[
\frac{\zeta^{(n+1)}(\nu)}{(n+1)!}A_{-1}
+\sum_{j=1}^{n}\frac{\zeta^{(j)}(\nu)}{j!}A_{n-j}.
\tag{CW1.13}
\] At \(1\), write \(\zeta(1+t)=t^{-1}+\sum_{j\ge0}c_jt^j\) and
\(a_f(1+t)=\sum_{k\ge1}a_kt^k\). The coefficient of \(t^n\) is \[
a_{n+1}+\sum_{j=0}^{n-1}c_ja_{n-j}.
\tag{CW1.14}
\] Empty sums are zero. Away from these exceptional points every
derivative is the full Leibniz sum
\(\sum_{j=0}^n\binom nj\zeta^{(j)}a_f^{(n-j)}\). Thus the comparison
retains derivative data as well as values.

Let \(\mathscr Z\) be the actual nontrivial zeros of the original
\(\zeta\), with their multiplicities \(m_\rho\), and let \[
\mathcal I=\{F\in\mathcal B:F^{(j)}(\rho)=0
\text{ for every }\rho\in\mathscr Z,\ 0\le j<m_\rho\},
\qquad \mathcal Q=\mathcal B/\mathcal I.
\tag{CW1.15}
\] CSP1.7 uses the proved SSI1--SSI7 exact-image theorem, independently
checked in ESI0--ESI8: \[
J=\overline J^{\,A},\qquad
\Sigma:S\xrightarrow{\sim}J\text{ topologically},\qquad
\Theta J=\mathcal I.
\tag{CW1.16}
\] The actual inverse in that theorem has, for \(F=\mathcal M_0a\), \[
f_F(x)=\frac1{2\pi}\int_{\mathbb R}
\frac{F(2+it)}{2\zeta(2+it)}x^{-2-it}\,dt,\quad x>0,\qquad
f_F(-x)=f_F(x),
\] \[
\frac{f_F^{(2r)}(0)}{(2r)!}
=\frac{F(-2r)}{2\zeta'(-2r)},\quad
f_F(0)=0,\quad \int_{\mathbb R}f_F=0.
\tag{CW1.17}
\] Its all-strip estimates and Schwartz remainder estimates are the
completed source theorem imported here, not a new assumption of closure.
Consequently \[
Q:=A/J=Q_{\rm alg}=Q_H
\xrightarrow[\ [a]\mapsto[\Theta a]\ ]{\sim}\mathcal Q
\tag{CW1.18}
\] is a topological isomorphism, and the former comparison kernel
\(\overline J/J\) is zero. The chain calculations below also make sense
algebraically before that theorem; continuity and the stated Fréchet
conclusions use the proved closed image.

\subsection{CW2. The actual sheaves, finite complexes, and
actions}\label{cw2.-the-actual-sheaves-finite-complexes-and-actions}

Let \(Y=\mathbb P^1(\mathbb C)\), \(D=\{0,\infty\}\),
\(U=\mathbb C^\times\), and \(j:U\hookrightarrow Y\),
\(i:D\hookrightarrow Y\). The map \(\pi:Y\to X=\{x_+,\eta,x_-\}\) sends
the poles to \(x_+,x_-\) and every other point to \(\eta\). Its
inverse-image sheaf is \[
\mathscr F=\pi^{-1}\Omega
=\underline A_Y\times_{i_*(A\oplus A)}i_*(V_+\oplus V_-).
\tag{CW2.1}
\] The fibre product uses the map
\((a,v_+,v_-)\mapsto(a|_0-r_+v_+,a|_\infty-r_-v_-)\). At the two poles
its stalks are \(V_\pm\); on \(U\) it is the constant sheaf with fibre
\(A\).

The actual subsheaf \(\mathscr F_J\) has the same pole stalks and
restrictions, with interior stalk \(J\). Put
\(\mathscr G=j_!\underline Q_U\). Quotienting the interior coefficients
and sending the pole stalks to zero gives the actual sheaf map \[
0\longrightarrow\mathscr F_J
\longrightarrow\mathscr F\xrightarrow{q_Y}\mathscr G
\longrightarrow0.
\tag{CW2.2}
\] Near a pole, the interior value of a section lies in
\(r_\pm(V_\pm)=J\). Its quotient is therefore zero near the pole,
exactly the extension-by-zero condition. The sequence is exact on every
stalk.

CSP2--CSP7 construct fine resolutions and their finite contractions. The
resulting finite complexes, with degrees \(0,1,2\), form the degreewise
exact row \[
\begin{array}{ccccccccc}
0&\to&V_+\oplus V_-&\xrightarrow{\mathrm{id}}&
V_+\oplus V_-&\to&0&\to&0\\
 &&\downarrow r_+-r_-&&\downarrow r_+-r_-&&\downarrow0\\
0&\to&J&\xrightarrow{\iota}&A&\xrightarrow{q}&Q&\to&0\\
 &&\downarrow0&&\downarrow0&&\downarrow0\\
0&\to&J&\xrightarrow{\iota}&A&\xrightarrow{q}&Q&\to&0.
\end{array}
\tag{CW2.3}
\] The three columns compute \(R\Gamma(Y,\mathscr F_J)\),
\(R\Gamma(Y,\mathscr F)\), and \(R\Gamma(Y,\mathscr G)\). The top row is
read as \(0\to D_J^0\to D_F^0\to D_G^0\to0\), with \(D_G^0=0\). All maps
are continuous.

For clarity, these are simultaneous models for the sheaf maps, not
independently chosen cohomology presentations. For any \(B=A,J,Q\), use
the fibre-product resolution with its specified pole spaces. It gives
the complex \[
B\oplus V_+\oplus V_-\longrightarrow B^2\longrightarrow B,
\qquad d^0(a,v_+,v_-)=(a-r_+v_+,a-r_-v_-),\quad d^1=0.
\] For \(B=Q\) the pole spaces are zero. The contraction \[
p^0(a,v_+,v_-)=(v_+,v_-),\quad
p^1(b_+,b_-)=b_--b_+,\quad p^2=\mathrm{id}
\] and the section \[
s^0(v_+,v_-)=(r_-v_-,v_+,v_-),\quad
s^1(b)=(-b,0),\quad s^2=\mathrm{id}
\tag{CW2.4}
\] commute with inclusion \(J\to A\) and quotient \(A\to Q\). With
\(h^1(b_+,b_-)=(b_-,0,0)\), direct substitution gives \(ps=1\) and
\(1-sp=dh+hd\). This proves the compatibility asserted in (CW2.3).

The coefficient actions are \[
T_a b(u)=b(u/a),\quad
\rho_+(a)(f,c_0,c_1)=(f(\cdot/a),c_0,ac_1),
\] \[
\rho_-(a)(h,d_0,d_1)=(a h(a\cdot),ad_0,d_1),\qquad a>0.
\tag{CW2.5}
\] They preserve every source and satisfy \(r_\pm\rho_\pm(a)=T_ar_\pm\).
The actual holomorphic covering \(b_n(z)=z^n\), for recovered positive
integers \(n\), has degree \(n\) on the sphere and multiplies an annular
angular period by \(n\). Combining that pullback with the coefficient
action gives \[
\mathsf B_n^0=\rho_+(n)\oplus\rho_-(n),\qquad
\mathsf B_n^1=T_n,\qquad
\mathsf B_n^2=nT_n.
\tag{CW2.6}
\] The same formula applies to the three columns where defined. Degree
two uses the positive complex orientation. In the notation \(B(-1)\),
the action is \(aT_a\); this records the factor just proved from
\(b_n\), without asserting an étale realization.

Integer operators are invertible on the computed cohomology. Their
inverses define the positive rational action consistently, and strong
continuity extends it to all \(a>0\). This is a cohomological action; no
nonintegral holomorphic map \(z\mapsto z^a\) is asserted.

\subsection{CW3. The connecting map of the fibre-product
presentation}\label{cw3.-the-connecting-map-of-the-fibre-product-presentation}

The other actual short exact sequence, CSP2.3, is \[
0\to\mathscr F\to
\underline A_Y\oplus i_*(V_+\oplus V_-)
\xrightarrow{\epsilon-r}i_*(A\oplus A)\to0.
\tag{CW3.1}
\] Constant-sheaf cohomology on the sphere is \(A,0,A\) in degrees
\(0,1,2\); a pole skyscraper has no positive cohomology. Hence its
complete nonzero long exact sequence is \[
0\to H^0(Y,\mathscr F)\to A\oplus V_+\oplus V_-
\xrightarrow d A^2
\xrightarrow{\delta_{\mathrm{fib}}}Q\to0,
\tag{CW3.2}
\] and the degree-two comparison is the identity \[
H^2(Y,\mathscr F)=A(-1)\xrightarrow{\mathrm{id}}
H^2(Y,\underline A_Y)=A(-1).
\tag{CW3.3}
\] Here \[
d(a,v_+,v_-)=(a-r_+v_+,a-r_-v_-),\qquad
\delta_{\mathrm{fib}}(b_+,b_-)=[b_--b_+].
\tag{CW3.4}
\] The connecting sign follows from the degree-one contraction \(p^1\)
in (CW2.4). Its kernel is exactly \[
\{(b_+,b_-):b_--b_+\in J\}.
\tag{CW3.5}
\] Indeed \(d\) has differences in \(J\). Conversely, for
\(b_--b_+=r_+v_+\), choose \(v_-=0\) and \(a=b_-\). Then
\(d(a,v_+,0)=(b_+,b_-)\). Every class \([b]\) is
\(\delta_{\mathrm{fib}}(0,b)\), so the image is all \(Q\).

This is an actual, generally nonzero connecting map. Its source action
is \(T_n\oplus T_n\) and its target action is \(T_n\), with no geometric
factor \(n\) in this degree. Its exact lifting quotient is \[
A^2/\operatorname{im}d\xrightarrow{\sim}Q,\qquad
[(b_+,b_-)]\mapsto[b_--b_+].
\tag{CW3.6}
\] Pairs in (CW3.5) lift through \(d\); other pairs do not. This
statement specifies this sheaf-presentation lifting problem, rather than
declaring a failure of a different global construction.

\subsection{CW4. The actual quotient-sheaf lift and its zero
boundary}\label{cw4.-the-actual-quotient-sheaf-lift-and-its-zero-boundary}

From (CW2.3), \[
H^0(Y,\mathscr F_J)=H^0(Y,\mathscr F)
=\{((\widehat h,c_0,c_1),(h,d_0,d_1))\},
\] \[
H^1(Y,\mathscr F_J)=0,\quad H^2(Y,\mathscr F_J)=J(-1),
\] \[
H^1(Y,\mathscr F)=Q,\quad H^2(Y,\mathscr F)=A(-1),
\] \[
H^0(Y,\mathscr G)=0,\quad H^1(Y,\mathscr G)=Q,\quad
H^2(Y,\mathscr G)=Q(-1),
\tag{CW4.1}
\] and all degree-three groups vanish. The actual long exact sequence is
therefore \[
0\to H^0(Y,\mathscr F_J)
\xrightarrow{\mathrm{id}}H^0(Y,\mathscr F)\to0,
\] \[
0\to H^1(Y,\mathscr F)=Q
\xrightarrow{\mathrm{id}}H^1(Y,\mathscr G)=Q
\xrightarrow{\delta^1}J(-1)
\xrightarrow{\iota}A(-1)
\xrightarrow{q}Q(-1)\xrightarrow{\delta^2}0.
\tag{CW4.2}
\] Every sign and map follows from the simultaneous complexes, including
\[
\boxed{\delta^1=0,\qquad \delta^2=0.}
\tag{CW4.3}
\]

Here is the cochain proof of the first equality, which does not assume a
section of \(q:A\to Q\). A class \(x\in H^1(Y,\mathscr G)=Q\) is
represented by the degree-one cocycle \(x\) of its zero-differential
complex. Choose any \(a\in A\) with \(q(a)=x\), which exists by the
definition of the quotient. Its differential in \(D_F\) is \(d_F^1a=0\).
Therefore the defining connecting cocycle in \(D_J^2\) is zero. Changing
\(a\) by an element of \(J\) still gives zero. The degree-one cohomology
class \([a]\in H^1(Y,\mathscr F)=A/J\) depends only on \(x\) and is
exactly \(x\).

Thus the actual lift is the continuous equivariant inverse of the
identity map in (CW4.2). It is unique at the cohomology-class level
because \(H^1(Y,\mathscr F_J)=0\). In degree two every \(x\in Q(-1)\)
also lifts, and its set of lifts is the nonempty affine fibre \[
\operatorname{Lift}_2(x)=q^{-1}(x)=a+J(-1).
\tag{CW4.4}
\] The equality describes all representatives, including their complete
kernel. It does not select a representative compatible with every
arithmetic action.

The degree-one obstruction receiver itself survives: \[
\operatorname{im}\delta^1=0,\qquad
\operatorname{coker}\bigl(H^1\mathscr F\to H^1\mathscr G\bigr)=0,
\qquad
\ker\bigl(H^2\mathscr F\to H^2\mathscr G\bigr)=J(-1).
\tag{CW4.5}
\] In the lifting formula of DC5, there is no additional invariant
quotient of \(H^1(Y,\mathscr G)\) in this sequence. Taking its actual
target as the classes to be lifted means that the further quotient is
the identity and its kernel is zero. The corresponding obstruction
quotient is therefore \(J(-1)/0\), while the obstruction map into it is
zero. This is the complete receiver; it has not been replaced by
\(Q(-1)\).

\subsection{CW5. Supported localization: every image, kernel, and
lift}\label{cw5.-supported-localization-every-image-kernel-and-lift}

For the two pole supports, the finite costalk complex is \[
L_D^0=V_+\oplus V_-,\quad L_D^1=A^2,\quad
L_D^2=A(-1)^2,\quad
d^0=(r_+,r_-),\quad d^1=0.
\] The support-to-global maps are \[
k^0=\mathrm{id},\qquad
k^1(b_+,b_-)=b_+-b_-,\qquad
k^2(c_+,c_-)=c_++c_-.
\tag{CW5.1}
\] The annular positive \(z\)-circle gives \[
R\Gamma(U,\underline A)=A[0]\oplus A(-1)[-1].
\] The local positive circle at infinity uses \(w=1/z\), so
\(d\arg w=-d\arg z\). With the positive complex local fundamental
classes, the actual localization boundaries are \[
\partial^0(a)=([a],[a]),\qquad
\partial^1(a)=(-a,a).
\tag{CW5.2}
\] These are precisely the cone maps of CSP7, including the negative
local boundary sign and the inverse-coordinate sign at infinity.

The resulting entire nonzero row is \[
\begin{aligned}
0\to E_+\oplus E_-&\to H^0(Y,\mathscr F)
\xrightarrow r A
\xrightarrow{\mathrm{diag}\,q}Q^2
\xrightarrow{\mathrm{diff}}Q\\
&\xrightarrow0 A(-1)
\xrightarrow{(-\mathrm{id},\mathrm{id})}A(-1)^2
\xrightarrow{\mathrm{sum}}A(-1)\to0.
\end{aligned}
\tag{CW5.3}
\] Here \(r\) sends a Fourier-graph section to its common restriction
and has image \(J\). Consequently:

\begin{enumerate}
\def\labelenumi{\arabic{enumi}.}
\tightlist
\item
  A degree-zero annular section \(a\in A\) extends globally exactly when
  \(a\in J\). Its obstruction is \(([a],[a])\), with image
  \(\operatorname{diag}Q\), kernel \(J\), and cokernel identified with
  \(Q\) by the difference map.
\item
  Every global degree-one class in \(Q\) lifts from the two supported
  groups \(Q^2\). The kernel is \(\operatorname{diag}Q\), and a
  continuous equivariant section is \[
  s_1(q)=\tfrac12(q,-q).
  \tag{CW5.4}
  \] The following restriction \(Q\to A(-1)\) is zero.
\item
  A degree-one annular class \(a\in A(-1)\) extends to global degree one
  exactly when \(a=0\). Its boundary is the injective map
  \(a\mapsto(-a,a)\), with full image the anti-diagonal. On that image
  its inverse is \((c_+,c_-)\mapsto(c_--c_+)/2\).
\item
  Every global degree-two class \(a\in A(-1)\) lifts from \(A(-1)^2\).
  The kernel is the anti-diagonal, and a continuous equivariant section
  is \[
  s_2(a)=\tfrac12(a,a).
  \tag{CW5.5}
  \]
\end{enumerate}

Each assertion is a direct kernel/image calculation in the displayed
row. In particular the quotient after the nonzero annular boundary is \[
A(-1)^2/\{(-a,a):a\in A(-1)\}
\xrightarrow{\sim}A(-1),\qquad[(b,c)]\mapsto b+c.
\tag{CW5.6}
\] It survives; the existence of the supported degree-two lift does not
assert that this receiver is zero.

For a single pole the supported degree-one map to \(Q\) is
\(+\mathrm{id}\) at \(0\), \(-\mathrm{id}\) at infinity, and the
supported degree-two map is \(+\mathrm{id}\) at both poles. These are
the maps of CSP7.7, so both single-support lifting maps are isomorphisms
in their displayed degrees.

Every scalar \(1/2\) in these sections belongs to the recovered
coefficient field. The sections do not divide the supporting datum.

The mirror signs are also retained. On \(Q\) both sphere swaps induce
\(-R\); on \(Q^2\) they induce \((q_+,q_-)\mapsto(Rq_-,Rq_+)\). Hence
\(s_1\) intertwines them. On \(A(-1)\), the holomorphic map
\(z\mapsto-1/z\) induces \(+R\), whereas the actual antiholomorphic map
\(z\mapsto-1/\overline z\) induces \(-R\). Their supported degree-two
maps are respectively \(+R\)-swap and \(-R\)-swap, so \(s_2\)
intertwines both. The annular actions are respectively \(-R,+R\),
exactly compatible with \(\partial^1(a)=(-a,a)\). These are coefficient
and orientation operations, not an exchanged label assigned to \(\tau\).

\subsection{CW6. Which character ranges are actually
separated}\label{cw6.-which-character-ranges-are-actually-separated}

Under \(\Theta\), coefficient dilation is multiplication by \(a^s\): \[
\Theta T_a=a^s\Theta,\qquad L F=sF.
\tag{CW6.1}
\] Fix \(a>1\) when discussing logarithmic-modulus weights, and define
\[
w_a(\lambda)=\frac{2\log|\lambda|}{\log a}\quad(\lambda\ne0).
\tag{CW6.2}
\] The actions themselves are defined for every \(a>0\), including the
identity at \(a=1\).

The full jet at an actual zero \(\rho\), in local coordinate
\(t=s-\rho\), has action \[
T_a=a^\rho\sum_{j=0}^{m_\rho-1}
\frac{(\log a)^j}{j!}M_t^j,\qquad M_t^{m_\rho}=0.
\tag{CW6.3}
\] On the normal quotient \(\mathcal Q(-1)\), the corresponding action
is \[
a^{\rho+1}\sum_{j=0}^{m_\rho-1}
\frac{(\log a)^j}{j!}M_t^j.
\tag{CW6.4}
\] The established open strip \(0<\Re\rho<1\) gives the disjoint ranges
\[
0<w_a(a^\rho)<2,\qquad
2<w_a(a^{\rho+1})<4.
\tag{CW6.5}
\] These preserve every nilpotent and do not impose \(\Re\rho=1/2\).
They are ranges of the actual zero characters, not a claim that an
infinite-dimensional complex representation is a mixed
finite-dimensional étale module.

The endpoint lines retain weights \(0,2\) on \(E_+\), and \(2,0\) on
\(E_-\). Their action is explicitly
\((c_0,c_1,d_0,d_1)\mapsto(c_0,ac_1,ad_0,d_1)\). The common-restriction
image of \(H^0(Y,\mathscr F)\) is \(J\), with action \(T_a\). Its
Fourier-graph parameter transforms as \(h(v)\mapsto a h(av)\).

It is essential to compute the characters of the actual degree-one
obstruction receiver \(J(-1)\), rather than assign to it (CW6.5). For
any \(\lambda\in\mathbb C\), let
\(m_\lambda=\operatorname{ord}_\lambda F_0\), equal to zero off
\(\mathscr Z\). For \(F\in\mathcal I\), local division defines \[
\ell_\lambda(F)=\left(\frac{F}{F_0}\right)(\lambda)
=\frac{F^{(m_\lambda)}(\lambda)}
{F_0^{(m_\lambda)}(\lambda)}.
\tag{CW6.6}
\] The denominator is nonzero. All lower numerator and denominator
derivatives vanish when \(m_\lambda>0\), giving the displayed ratio.
Evaluation of a fixed derivative is continuous for the strip topology by
the Cauchy integral formula on a small fixed circle. Thus
\(\ell_\lambda\) is continuous. Moreover \[
\ell_\lambda(F_0)=1,\qquad
\ell_\lambda(a^sF)=a^\lambda\ell_\lambda(F).
\tag{CW6.7}
\] The first equality proves surjectivity to \(\mathbb C\); the second
follows either from local division or from Leibniz's formula with all
lower vanishing derivatives retained.

Consequently \(\mathcal I(-1)\), and hence \(J(-1)\), has a continuous
one-dimensional quotient character \[
a\longmapsto a^{\lambda+1}
\quad\text{for every }\lambda\in\mathbb C.
\tag{CW6.8}
\] At \(0,1,-2r\) the order \(m_\lambda\) is zero and the nonzero
denominators are exactly (CW1.11), so there is no exceptional gap in
this statement. For an actual original zero \(\rho\), the choice
\(\lambda=\rho-1\) gives the quotient character \(a^\rho\), the same
character that occurs in \(\mathcal Q\). The point \(\rho-1\) is an
evaluation point of the full original function; it has not been asserted
to be a zero of \(\zeta\).

The same elementary evaluation on \(\mathcal B\) gives all characters
\(a^\lambda\) as quotient characters of \(A\), and all \(a^{\lambda+1}\)
as quotient characters of \(A(-1)\). Therefore the actual targets
\(J(-1)\) and \(A(-1)\) in (CW4.2) and (CW5.3) do not satisfy a bound
placing all their continuous character quotients above the original-zero
range. The separated space is the \emph{quotient} \(Q(-1)\) of the full
normal source, with its complete original-zero ideal removed.

This calculation concerns quotient characters. It does not claim that
\(J\) or \(A\) has corresponding eigenvectors. In fact they have no
nonzero arithmetic eigenvectors at all: an eigenfunction
\(F\in\mathcal B\) for one \(a>1\) would satisfy \((a^s-c)F(s)=0\). The
first factor is a nonzero entire function, so the product identity
forces \(F=0\). The same argument, with a power of \(a^s-c\), excludes
nonzero generalized eigenvectors in these entire-function spaces.

\subsection{CW7. A global vanishing theorem for the actual boundary
maps}\label{cw7.-a-global-vanishing-theorem-for-the-actual-boundary-maps}

There is a complete global vanishing calculation even though the actual
targets in CW6 do not have the proposed separated quotient-character
bound. We prove it using the full quotient topology and a resolvent,
without assuming that finite zero-jet subspaces are dense.

For \(\mu\notin\mathscr Z\), define on \(\mathcal B\) \[
\mathcal R_\mu F(s)=
\frac{F(s)-F_0(s)F(\mu)/F_0(\mu)}{\mu-s}.
\tag{CW7.1}
\] The numerator vanishes at \(s=\mu\), so this is entire. It belongs to
\(\mathcal B\) continuously. To verify the topology directly, put
\(H=F-F_0F(\mu)/F_0(\mu)\). Outside \(|s-\mu|<1\), division by \(\mu-s\)
is bounded by one. Inside that disk, \[
\frac{H(s)}{\mu-s}
=-\int_0^1H'(\mu+t(s-\mu))\,dt.
\] Cauchy's estimate bounds \(|H'|\) there by
\(\sup_{|z-\mu|\le2}|H(z)|\). For \(M'=\max(M,|\Re\mu|+2)\), this gives
\[
b_{M,N}(\mathcal R_\mu F)
\le b_{M,N}(H)+(2+|\Im\mu|)^N b_{M',0}(H).
\tag{CW7.2}
\] The scalar \(F(\mu)/F_0(\mu)\) is bounded by
\(b_{M',0}(F)/|F_0(\mu)|\). Substituting it in (CW7.2) gives a finite
sum of continuous seminorm bounds on \(F\), with the complete \(F_0\)
retained.

If \(F\in\mathcal I\), its numerator is in \(\mathcal I\), and division
by \(\mu-s\) preserves all original-zero vanishing orders. Thus (CW7.1)
descends continuously to \(\mathcal Q\). Direct multiplication gives \[
(\mu-L)\mathcal R_\mu[F]=[F],\qquad
\mathcal R_\mu(\mu-L)[F]=[F].
\tag{CW7.3}
\] The discarded term in the first equality is explicitly
\(F_0F(\mu)/F_0(\mu)\in\mathcal I\); no factor of the original function
has been suppressed. This rederives the needed RZ inverse on its exact
full domain.

Let \(H:\mathcal Q\to\mathcal B(-1)\) be continuous and complex-linear
and commute with the recovered arithmetic actions. If equivariance is
initially required only for positive integers, invertibility gives it
for positive rationals; continuity and strong continuity of \(T_a\) give
it for every \(a>0\). Differentiation at \(a=e^t,t=0\) is justified by
\[
|e^{ts}-1-ts|\le\tfrac12t^2|s|^2e^{|t||\Re s|},
\] which proves convergence in every strip seminorm, and then in every
quotient seminorm. The target generator is \(L+1\). Consequently \[
H L=(L+1)H.
\tag{CW7.4}
\] For a fixed \(\lambda\in\mathbb C\), evaluation gives \[
\operatorname{ev}_\lambda H\bigl(L-(\lambda+1)\bigr)=0.
\tag{CW7.5}
\] When \(\lambda+1\notin\mathscr Z\), (CW7.3) makes the operator in
parentheses onto \(\mathcal Q\). Hence \(H(x)(\lambda)=0\) for every
\(x\in\mathcal Q\). For each \(x\), the entire function \(H(x)\)
therefore vanishes on the open complement of the discrete set
\(\mathscr Z-1\), and is identically zero. We have proved \[
\boxed{\operatorname{Hom}_{\mathrm{cont},T}
(\mathcal Q,\mathcal B(-1))=0.}
\tag{CW7.6}
\] Since \(\mathcal I(-1)\) is an invariant subspace, inclusion gives \[
\boxed{\operatorname{Hom}_{\mathrm{cont},T}
(\mathcal Q,\mathcal I(-1))=0.}
\tag{CW7.7}
\] Transport through \(\Theta\) proves the corresponding statements for
\(Q\to A(-1)\) and \(Q\to J(-1)\). Thus both the actual restriction
\(H^1(Y,\mathscr F)\to H^1(U,\underline A)\) and the actual connecting
map \(\delta^1:H^1(Y,\mathscr G)\to H^2(Y,\mathscr F_J)\) also vanish by
this global operator calculation. Their explicit cochain proofs in
CW4--CW5 identify which maps are being tested.

No assertion about a direct sum or unrestricted product of zero blocks
entered the proof. Nor did overlapping quotient characters produce a
nonzero map: the distinction between a quotient character and an
entire-function eigenvector is retained.

For comparison, the separated original-zero targets also admit a global
proof: \[
\operatorname{Hom}_{\mathrm{cont},T}(\mathcal Q,\mathcal Q(-1))=0.
\tag{CW7.8}
\] Indeed, for a continuous complex-linear equivariant \(H\), take the
full \(m_\rho\)-jet of its target at an actual \(\rho\). The target
generator there is \((\rho+1)I+M_t\), so that jet kills
\((L_{\mathrm{target}}-(\rho+1))^{m_\rho}\). Equivariance and (CW7.3),
with \(\mu=\rho+1\notin\mathscr Z\), imply that every target jet of
\(H(x)\) is zero. The joint original-zero jets are faithful on
\(\mathcal Q\) by (CW1.15). Thus \(H=0\). This proves the global
assertion behind the separated ranges in (CW6.5), including every
multiplicity.

\subsection{CW8. The normal quotient lifts classes but has no
equivariant linear
section}\label{cw8.-the-normal-quotient-lifts-classes-but-has-no-equivariant-linear-section}

The actual degree-two row is \[
0\to J(-1)\to A(-1)\xrightarrow{q}Q(-1)\to0,
\quad\text{or, under }\Theta,\quad
0\to\mathcal I(-1)\to\mathcal B(-1)
\xrightarrow{q}\mathcal Q(-1)\to0.
\tag{CW8.1}
\] Every class lifts by its definition as a quotient; the whole fibre is
(CW4.4). We now calculate the distinct question of a simultaneous
arithmetic-equivariant linear choice of representatives.

For any actual zero \(\rho\), put \[
e_\rho(s)=F_0(s)/(s-\rho).
\tag{CW8.2}
\] It is entire and in \(\mathcal B\), by the same division estimate
near \(\rho\) used above and polynomial division away from \(\rho\). Its
order at \(\rho\) is \(m_\rho-1\), so \([e_\rho]\ne0\) in
\(\mathcal Q\). Moreover \[
(L-\rho)[e_\rho]=[F_0]=0,\qquad
T_a[e_\rho]=a^\rho[e_\rho].
\tag{CW8.3}
\] For the second equality, \((a^s-a^\rho)/(s-\rho)\) is entire and a
multiplier with polynomial strip bounds away from a fixed disk;
multiplied by \(F_0\), it lies in \(\mathcal I\), with all required
orders at every original zero.

The actual zero set is nonempty without an RH assumption. One way using
the retained source proof is S2's order-at-most-one Hadamard
factorization and \(F_0(s)=F_0(1-s)\). If it were zero-free, that
factorization would give \(F_0=e^{a+bs}\); the symmetry forces \(b=0\).
This contradicts \(F_0(0)=1/8\) and \(F_0(2)=\pi/24\), since
\(\pi\ne3\). This argument uses the original full function and does not
assign a location to a fictional zero.

An arithmetic-equivariant linear section of (CW8.1) would send the
nonzero class \([e_\rho]\), with its twisted eigencharacter
\(a^{\rho+1}\), to a nonzero entire function \(F\). For one \(a>1\),
equivariance would require \[
a^{s+1}F(s)=a^{\rho+1}F(s).
\tag{CW8.4}
\] The nonzero entire factor \(a^{s+1}-a^{\rho+1}\) forces \(F=0\), a
contradiction. Thus \[
\boxed{\text{The actual normal quotient has no arithmetic-equivariant
linear section.}}
\tag{CW8.5}
\] Continuity is not needed for this no-section statement. In the
continuous category, twisting both source and target by the same
character and applying the proof of CW7 with target generator \(L\) also
gives \[
\operatorname{Hom}_{\mathrm{cont},T}
(\mathcal Q(-1),\mathcal B(-1))=0.
\tag{CW8.6}
\]

This is a statement about simultaneous compatible representatives, not
about the existence of individual lifts. The latter are all present,
with the exact affine fibres \(a+J(-1)\). In particular (CW8.5) does not
contradict the degree-one identity lift in CW4, the surjective
degree-two map, or Deligne's surjectivity theorem. DC8 explicitly
distinguished its canonical section into the intermediate \(B_i\) from
an unspecified choice of lifts into the special-fibre group \(A_i\).

\subsection{CW9. The precise separated receiver obtained from the actual
normal
quotient}\label{cw9.-the-precise-separated-receiver-obtained-from-the-actual-normal-quotient}

The geometric normal quotient is \(\mathcal Q(-1)\), with action
(CW6.4). It has the exact map into the GSL extension, rather than an
identification with the kernel \(J(-1)\) of (CW8.1).

Let \[
\mathcal I_+=\{F\in\mathcal B:
F^{(j)}(\rho+1)=0\quad(0\le j<m_\rho)\},\qquad
\mathcal E_+=\mathcal B/(\mathcal I\cap\mathcal I_+).
\tag{CW9.1}
\] GSL4--GSL6 construct the entire multiplier \(E_+\), with jet one on
the original divisor and jet zero on the divisor translated by \(+1\),
and prove its strip bounds \[
|E_+(s)|+|1-E_+(s)|
\le C_M(1+|\Im s|)^{d_M}\quad(|\Re s|\le M).
\tag{CW9.2}
\] These give continuous multiplication on \(\mathcal B\):
\(b_{M,N}(E_+F)\le C_M b_{M,N+d_M}(F)\), and the analogous bound for
\(1-E_+\). Those completed global bounds, not a formal blockwise
interpolation, are the input from GSL.

With \(VF(s)=F(s-1)\), the complete maps are \[
j_{\mathrm{high}}:\mathcal Q(-1)\to\mathcal E_+,\quad
[h]\mapsto[(1-E_+(s))h(s-1)],
\] \[
\pi_{\mathrm{high}}[F]=[F]_{\mathcal I},\qquad
s_{\mathrm{high}}[f]=[E_+f].
\tag{CW9.3}
\] They give \[
0\to\mathcal Q(-1)\xrightarrow{j_{\mathrm{high}}}
\mathcal E_+\xrightarrow{\pi_{\mathrm{high}}}\mathcal Q\to0,
\qquad \pi_{\mathrm{high}}s_{\mathrm{high}}=1.
\tag{CW9.4}
\] Independence of representatives follows from the two full jet
identities, including all derivatives. The inverse comparison is \[
\mathcal E_+\longrightarrow\mathcal Q\oplus\mathcal Q(-1),
\qquad
[F]\longmapsto([F]_{\mathcal I},[F(s+1)]_{\mathcal I}),
\tag{CW9.5}
\] whose inverse is \[
([f],[h])\longmapsto
[E_+(s)f(s)+(1-E_+(s))h(s-1)].
\tag{CW9.6}
\] The two inverse identities hold at every original and translated jet.
Continuity follows from (CW9.2), translation
\(b_{M,N}(VF)\le b_{M+1,N}(F)\), and the quotient seminorms.
Equivariance retains the factor \[
a\,a^{s-1}=a^s.
\tag{CW9.7}
\]

The complete original product defining the two divisors is \[
F_0(s)F_0(s-1)=
\left[\frac{s(s-1)}8\pi^{-s/2}\Gamma(s/2)\zeta(s)\right]
\left[\frac{(s-1)(s-2)}8\pi^{-(s-1)/2}
\Gamma((s-1)/2)\zeta(s-1)\right].
\tag{CW9.8}
\] The translated factor has values \(1/8\) at \(s=1,2\), value (CW1.11)
at \(s=1-2r\), and the same value at \(s=2+2r\). Its derivatives are
exactly \(F_0^{(j)}(s-1)\); derivatives of the product are the full sum
\[
\sum_{j=0}^{n}\binom njF_0^{(j)}(s)F_0^{(n-j)}(s-1).
\tag{CW9.9}
\] Original and translated endpoint and trivial-zero contributions
therefore remain explicit.

The actual geometric normal map is \[
\mathfrak n:A(-1)\to\mathcal E_+,\qquad
a\mapsto[(1-E_+(s))(\Theta a)(s-1)].
\tag{CW9.10}
\] Its factorization is exactly \[
A(-1)\xrightarrow{q}\mathcal Q(-1)
\xrightarrow{j_{\mathrm{high}}}\mathcal E_+.
\tag{CW9.11}
\] The second arrow is injective, so \[
\ker\mathfrak n=J(-1),\qquad
\operatorname{im}\mathfrak n=\ker\pi_{\mathrm{high}}.
\tag{CW9.12}
\] Thus the separated GSL kernel is the complete quotient of the actual
normal source by its full original summation image. It is not the
summation image itself.

The comparison on the actual connecting-map receiver can be computed
without discarding that receiver. If
\(\iota:J(-1)\hookrightarrow A(-1)\) is the actual map in (CW4.2), then
\[
q\iota=0,\qquad
\mathfrak n\iota=j_{\mathrm{high}}q\iota=0,\qquad
\mathfrak n\iota\delta^1=0.
\tag{CW9.15}
\] The middle map is zero on the entire receiver, before considering the
value of \(\delta^1\). Therefore vanishing after this comparison is not
an injective test for vanishing in \(J(-1)\). The actual vanishing of
\(\delta^1\) is instead proved by (CW4.3) and independently by (CW7.7).
This identifies the exact comparison map and its kernel, rather than
treating the two receiver presentations as unrelated.

The actual two-degree groups give the continuous equivariant isomorphism
\[
H^1(Y,\mathscr F)\oplus H^2(Y,\mathscr G)
\xrightarrow{\sim}\mathcal E_+
\tag{CW9.13}
\] with formula (CW9.6). This map forgets the distinction between the
two cohomological degrees only at its displayed direct-sum target. The
degree-one connecting morphism of the actual sheaf sequence remains
\(\delta^1:Q\to J(-1)\), not a new morphism \(Q\to Q(-1)\).

For precision, the split algebraic cross on this receiver is \[
0\to\mathcal Q(-1)\to\mathcal Q\oplus\mathcal Q(-1)
\xrightarrow{\mathrm{pr}_1}\mathcal Q\to0,\qquad
\partial_{\mathrm{sum}}=\mathrm{pr}_2.
\tag{CW9.14}
\] The actual degree-one group enters by \(x\mapsto(x,0)\), and its
boundary is zero. The image of the kernel under
\(\partial_{\mathrm{sum}}\) is all \(\mathcal Q(-1)\), so the DC-style
obstruction quotient of this particular cross is zero. Transport through
(CW9.13) gives exactly (CW9.4), whose continuous complex-linear
equivariant section is unique by (CW7.8).

This cross is constructed from the computed cohomology groups and the
exact normal quotient. It is not substituted for (CW4.2) or (CW5.3). In
particular, the separated kernel in (CW9.14) is visible only after the
stated quotient and assembly of two degrees.

\subsection{CW10. Exact comparison with Deligne's invariant-cycle
mechanism}\label{cw10.-exact-comparison-with-delignes-invariant-cycle-mechanism}

DC3--DC9 derive, in Deligne's arithmetic reduction, the two linked rows
\[
0\to K_i=H^{i-1}(X_{\bar\eta},E)_I(-1)
\xrightarrow{j_i}B_i=H^i(X_\eta,E)
\xrightarrow{\pi_i}C_i=H^i(X_{\bar\eta},E)^I\to0,
\] \[
A_i=H^i(X_s,E)\xrightarrow{\alpha_i}B_i
\xrightarrow{\partial_i}
O_i=H^{2N-i-1}(X_s,E)^\vee(-N).
\tag{CW10.1}
\] Here \(N\) is the dimension of the smooth total model, not the
generic fibre dimension. The invariant-cycle theorem's hypotheses are
those of DC1--DC2, including the curve-trait setting and arithmetic
reduction; they are not transferred to the sphere sheaf by its notation.

Its exact lifting obstruction is \[
\operatorname{obs}_i(c)=[\partial_i b]
\in O_i/\partial_i j_i(K_i),\qquad \pi_i(b)=c.
\tag{CW10.2}
\] Changing \(b\) by \(j_i(k)\) changes the boundary by
\(\partial_i j_i(k)\), proving well-definedness. Conversely,
\(\operatorname{obs}_i(c)=0\) permits subtracting such a \(j_i(k)\) and
then lifting through \(\alpha_i\) by localization. Thus \[
\operatorname{coker}(\pi_i\alpha_i)
\simeq\operatorname{im}\partial_i/\partial_i j_i(K_i).
\tag{CW10.3}
\] The geometric theorems prove \[
C_i,\ A_i\text{ have weights }\le i,\qquad
K_i,\ O_i\text{ have weights }\ge i+1.
\tag{CW10.4}
\] The bound on \(O_i\) is the full calculation \[
-(2N-i-1)+2N=i+1;
\] the bound on \(K_i\) comes from dual invariant weights followed by
the retained \((-1)\) twist. Exactness of the weight cutoff in the
finite-dimensional mixed category gives \[
W_iB_i\xrightarrow{\sim}C_i,\quad
\partial_i(W_iB_i)=0,\quad
\operatorname{im}\alpha_i=W_iB_i,\quad
\partial_i j_i(K_i)=\operatorname{im}\partial_i.
\tag{CW10.5}
\] The whole receiver need not vanish: \[
O_i/\partial_i j_i(K_i)\simeq\ker\operatorname{sp}_{i+1}.
\tag{CW10.6}
\]

The actual CSP sequence now has a fully determined comparison:

{\def\LTcaptype{none} % do not increment counter
\begin{longtable}[]{@{}
  >{\raggedright\arraybackslash}p{(\linewidth - 6\tabcolsep) * \real{0.2500}}
  >{\raggedright\arraybackslash}p{(\linewidth - 6\tabcolsep) * \real{0.2500}}
  >{\raggedright\arraybackslash}p{(\linewidth - 6\tabcolsep) * \real{0.2500}}
  >{\raggedright\arraybackslash}p{(\linewidth - 6\tabcolsep) * \real{0.2500}}@{}}
\toprule\noalign{}
\begin{minipage}[b]{\linewidth}\raggedright
Actual calculation
\end{minipage} & \begin{minipage}[b]{\linewidth}\raggedright
Classes being lifted
\end{minipage} & \begin{minipage}[b]{\linewidth}\raggedright
Actual receiver or kernel
\end{minipage} & \begin{minipage}[b]{\linewidth}\raggedright
Calculated result
\end{minipage} \\
\midrule\noalign{}
\endhead
\bottomrule\noalign{}
\endlastfoot
Quotient-sheaf degree one, (CW4.2) & \(H^1\mathscr G=Q\) &
\(H^2\mathscr F_J=J(-1)\) & Boundary zero; unique cohomology lift, the
identity \\
Quotient-sheaf degree two, (CW8.1) & \(H^2\mathscr G=Q(-1)\) & Kernel
\(J(-1)\); next boundary target \(0\) & Every class lifts; all lifts
form \(a+J(-1)\) \\
Support to global degree one, (CW5.3) & \(H^1\mathscr F=Q\) & Following
receiver \(A(-1)\) & Following map zero; section
\(q\mapsto(q/2,-q/2)\) \\
Global to annular degree zero & \(A\) & \(Q^2\) & Boundary
\(\operatorname{diag}q\), image \(\operatorname{diag}Q\) \\
Global to annular degree one & \(A(-1)\) & \(A(-1)^2\) & Injective
boundary \(a\mapsto(-a,a)\) \\
Normal quotient assembled with degree one, (CW9.14) & \(Q\) & \(Q(-1)\)
& Separated continuous split receiver; obstruction quotient zero \\
\end{longtable}
}

The first and third zero maps are proved both by their actual complexes
and by the global continuous intertwiner calculation of CW7. Their full
targets \(J(-1)\) and \(A(-1)\) have the character quotients of CW6, so
Deligne's finite-dimensional weight-cutoff proof (CW10.4) is not the
proof of those target bounds here. This is an exact determination of the
present objects and maps, not an assertion that no further geometric
comparison can exist.

There is already a precise morphism to the separated normal quotient:
(CW9.10)--(CW9.12), with the entire kernel \(J(-1)\) and the entire
image \(\ker\pi_{\mathrm{high}}\) calculated. It is this quotient and
the two-degree comparison, rather than a renaming of \(J(-1)\), that
relates the actual CSP geometry to GSL's separated row.

No map in this calculation forces the original characters \(a^\rho\) all
to have weight one. The surviving original quotient still has its
complete original zero jets, their nilpotent actions, and their actual
strip positions. The results prove the stated actual lifts and their
exact receivers; they do not assume a critical-line conclusion or
construct a violation of it.

\subsection{CW11. Prerequisites, source use, and retained
scope}\label{cw11.-prerequisites-source-use-and-retained-scope}

The operations used in this derivation have the following constructed
prerequisites:

{\def\LTcaptype{none} % do not increment counter
\begin{longtable}[]{@{}
  >{\raggedright\arraybackslash}p{(\linewidth - 2\tabcolsep) * \real{0.5000}}
  >{\raggedright\arraybackslash}p{(\linewidth - 2\tabcolsep) * \real{0.5000}}@{}}
\toprule\noalign{}
\begin{minipage}[b]{\linewidth}\raggedright
Operation
\end{minipage} & \begin{minipage}[b]{\linewidth}\raggedright
Exact prerequisite and proof
\end{minipage} \\
\midrule\noalign{}
\endhead
\bottomrule\noalign{}
\endlastfoot
Vector sums, differences and \(1/2\) & Recovered complex coefficient
spaces \(A,S,V_\pm,Q\); CW1--CW2. None acts on \(\tau\). \\
Sheaf kernel and quotient & Stalkwise exact fibre-product and
extension-by-zero maps; CSP2/CSP10 and CW2.1--CW2.3. \\
Long exact sequences and boundaries & Simultaneous finite complexes and
their compatible contractions; CW2.3--CW2.4, CW3--CW5. \\
Ordinary equals Hausdorff cohomology & Actual closed-image/inverse
theorem SSI1--SSI7 and ESI0--ESI8, imported through CSP1.7 with its full
inverse CW1.17. \\
Normal degree character & Actual positive-degree covering
\(z\mapsto z^n\), annular period and oriented sphere integration; CSP8
and CW2.6. \\
Original-zero weights & Exact raw Mellin comparison, original zero ideal
and full nilpotent action; CW1.6--CW1.18 and CW6. \\
Evaluation on the full kernel & Local division by the retained full
\(F_0\), with all exceptional values; CW6.6--CW6.8. \\
Global vanishing of continuous intertwiners & Continuous full resolvent,
strong differentiability and holomorphic uniqueness; CW7. No density of
finite blocks is used. \\
Shifted normal comparison & GSL's constructed multiplier with strip
bounds, full translated factor and continuous CRT inverse;
CW9.1--CW9.12. \\
Deligne comparison & The actual invariant/coinvariant, duality and
weight maps of DC3--DC9; full total-dimension and Tate factors retained
in CW10. \\
\end{longtable}
}

The new statements about absence of equivariant linear sections are
restricted to the actual coefficient row (CW8.1). They preserve the
whole original arithmetic and do not replace it by a different timing
history. They do not deny individual lifting, and they do not identify
that coefficient-row section problem with an unconstructed obstruction
in the user's complete construction.

The concrete mathematical conclusions are the actual cohomological
identity lift, both zero connecting maps of (CW4.2), the full
localization maps and their explicit sections, the global vanishing
(CW7.6)--(CW7.8), the exact normal lift fibres, and the quotient
morphism to the separated receiver. Every comparison keeps the support,
separate chart labels, full source kernel, original zeta factors,
multiplicities, and cohomological degrees at the points where they are
used.
